\documentclass[11pt,a4paper]{article}
\usepackage[margin=27mm]{geometry}
\usepackage{amsmath,amssymb,amsthm,hyperref}
\newtheorem{lemma}{Lemma}\newtheorem{theorem}[lemma]{Theorem}
\newcommand{\Sz}{\operatorname{Sz}}
\newcommand{\gen}[1]{\langle#1\rangle}
\title{Obstructions to five prescribed Suzuki constructions}
\author{Human Compression project}\date{}
\begin{document}\maketitle
\begin{abstract}
The higher-rank constructions in the fixed Suzuki pregeometry article
prescribe a full intersection pattern, not just abstract parabolic types.
We separate an elementary defect in three printed intersection formulas
from the nonexistence of the intended patterns. Two previously established
normalizer arguments exclude those patterns and the related first and
sixth constructions. The arguments are reproduced here so that the source
disposition does not depend on a label in a working ledger.
\end{abstract}

Write $q=2^n$, with $n$ odd, and work in the classical Suzuki matrix model.
For odd $d\mid n$, let $S_d$ be the standard subgroup obtained by restricting
parameters to $\mathbb F_{2^d}$. For $d>1$ this is the simple group
$\Sz(2^d)$; $S_1=\Sz(2)$ has order 20 and is isomorphic to $C_5:C_4$.
Let $H_q$ be the split torus and $w$ its standard involutory normalizer.
Put $\Delta_d=H_{2^d}\gen w$, of order $2(2^d-1)$, where $H_{2^d}$ denotes
the unique subgroup of $H_q$ of order $2^d-1$.

We retain the classical conjugacy of subfield groups, conjugacy of their
maximal split dihedrals, the split-torus centralizer and normalizer, and
the involution and Sylow-two structure. These are explicit earlier inputs,
not consequences of a finite sample. In the standard coordinates the
unique Bruhat parameters give
\begin{equation}\label{field}
 S_a\cap S_b=S_{\gcd(a,b)}.
\end{equation}
Indeed a matrix belongs to both groups precisely when its uniquely
recovered parameters belong to both fields. This also proves the statement
when the common field is $\mathbb F_2$.

\section{Two rigidity arguments}
\begin{lemma}\label{split}
For $d>1$, the normalizer of $\Delta_d$ in $\Sz(q)$ is $\Delta_d$.
Exactly one conjugate of $S_d$ contains $\Delta_d$.
\end{lemma}
\begin{proof}
The odd cyclic subgroup of $\Delta_d$ is characteristic, and its centralizer
is $H_q$. Every normalizer of $\Delta_d$ therefore belongs to
$H_q\gen w$. For $h\in H_q$, conjugation preserves $\Delta_d$ if and only
if $h^2\in H_{2^d}$. The quotient $H_q/H_{2^d}$ has odd order,
so this is equivalent to $h\in H_{2^d}$. This proves the normalizer claim.
If a conjugate of $S_d$ contains $\Delta_d$, pull the latter back into $S_d$.
It is a maximal split dihedral there and can be moved to the standard
one by internal conjugacy. After this adjustment the defining conjugator
normalizes $\Delta_d$, hence belongs to $\Delta_d\leq S_d$. The conjugate is $S_d$.
\end{proof}

\begin{theorem}\label{alignment}
Suppose $a,b>1$ are odd divisors of $n$. Let $S,T$ be conjugates of
$S_a,S_b$. Suppose one split dihedral subgroup contains maximal split
dihedrals of $S$ and $T$ with a common involution. Then $S,T$ are
simultaneously conjugate to $S_a,S_b$. In particular,
$S\cap T\cong\Sz(2^{\gcd(a,b)})$.
\end{theorem}
\begin{proof}
Conjugate the ambient split-torus normalizer to $H_q\gen w$.
Conjugating further by $H_q$ moves the common reflection to $w$,
since squaring is bijective in this odd cyclic torus. Its rotation
subgroups of orders $2^a-1$ and $2^b-1$ are unique. The two prescribed
dihedrals are therefore $\Delta_a,\Delta_b$. Lemma~\ref{split} forces the two
subfield groups to be standard. Now apply \eqref{field}.
\end{proof}

For the second argument use the standard Sylow-two subgroup $U_q$ and
its center $Z_q=\{u(0,t):t\in\mathbb F_q\}$. Its nonidentity central
elements are precisely its involutions. Fix $z=u(0,1)$.
The retained involution-centralizer theorem says $C_G(z)=U_q$.

\begin{lemma}\label{center}
If a conjugate $K$ of $S_d$ contains $z$, then the identity together
with the involutions in $C_K(z)$ is exactly $Z_{2^d}$.
Consequently, conjugates of $S_a,S_b$ with $\gcd(a,b)=1$ cannot
have a Klein four subgroup in their intersection.
\end{lemma}
\begin{proof}
The involutions of $S_d$ form one conjugacy class. Adjusting a defining
conjugator internally therefore writes $K=S_d^v$ with
$v\in C_G(z)=U_q$. Thus $C_K(z)=U_{2^d}^v$. Its identity and
involutions form $Z_{2^d}^v=Z_{2^d}$, because $U_q$ centralizes $Z_q$.
If two subfield groups share an involution, first conjugate it to $z$.
Their common commuting involutions lie in
$Z_{2^a}\cap Z_{2^b}=Z_{2^{\gcd(a,b)}}$.
For coprime degrees this group has order two. A Klein four group would
supply three nonidentity commuting involutions, a contradiction.
\end{proof}

\section{The rank-three patterns and their extensions}
The fixed article L110 reproduces as R4, R5 and R6 three patterns in
$\Sz(2^{ab})$, where $a,b$ are distinct odd primes. They require
two subfield groups $S\cong\Sz(2^a)$ and $T\cong\Sz(2^b)$ and a
split dihedral $D$ with
\[
 |D|=2(2^a-1)(2^b-1),\quad
 S\cap D\cong D_{2(2^a-1)},\quad
 T\cap D\cong D_{2(2^b-1)},\quad
 S\cap T\cap D\cong C_2.
\]
They prescribe $S\cap T$ to be, respectively, $D_{10}$,
$C_2^2$, or $C_4$. Theorem~\ref{alignment} instead gives
$S\cap T\cong S_1$, of order 20. None of these three full patterns exists.
The primality conditions on $2^a-1$ and $2^b-1$ are not needed for this
obstruction.

The pattern R8 has the same coprime subfield degrees and prescribes
$S\cap T\cong C_4\times C_2$. This group contains a Klein four group,
so Lemma~\ref{center} excludes it. No primality condition on the
chosen nonsplit-torus factors is needed for that conclusion.

The implication for a prescribed extension is elementary but must be
stated with its hypotheses. If a larger family of actual subgroups has a
residue with one of these full intersection patterns, intersecting the
family with the residue stabilizer would realize the forbidden pattern.
Hence such a larger family cannot exist. This reasoning concerns the
full prescribed lattice; a bare list of abstract subgroup types does not
specify the required intersections.

L110 Section 3.3.1, printed pp.12--13, derives Constructions 2, 3, 4 and 6
from residues R4, R5, R6 and R8, respectively. The introduction to the
constructions on p.15 explicitly requires the full sublattice obtained
from the corresponding rank-four pattern. In higher rank, fixing all
but the appropriate two subfield types in these prescribed patterns
leaves the corresponding rank-three pattern in a subfield stabilizer.
The obstruction therefore applies to every intended full lattice in
these four constructions, not just to a convenient coordinate placement.

Construction 1 requires a separate check. The rank-four lattice at the
end of Section 3.1, printed p.11, has two subfield groups of distinct
prime degrees and the split dihedral of order
$2(2^a-1)(2^b-1)$. Their intersections with that dihedral are their
maximal split dihedrals, with a common involution. The two subfield
groups are asserted to intersect in $D_{10}$ (the case $n=5$ in that
paragraph). This is exactly the contradiction in
Theorem~\ref{alignment}; the additional affine parabolic cannot repair it.
Construction 1 extends this prescribed rank-four lattice. For more than
two subfield types, fixing the other subfield types leaves the same
rank-four residue. Thus none of its intended full lattices exists either.
Figure 1 is explicitly conditional on existence and flag transitivity;
the impossible lattice supplies no realized diagram.

\section{The literal formulas are a separate issue}
Let the odd field degree be the squarefree product $p_1\cdots p_m$
and put $n_i=2^{p_i}-1$, as in Constructions 2--4. Their printed
intersection runs from $j=0$ through $m$, so includes
$G_m\cong D_{2\prod_i n_i}$. The printed values are $D_{10}$,
$D_4=C_2^2$, and $C_4$, respectively. All three are impossible even
before any subfield argument. Since each $p_i$ is odd,
$5\nmid2^{p_i}-1$; hence the ambient dihedral order is not divisible
by 5. Its two-part is exactly 2, so it has no subgroup of order 4.
Lagrange's theorem excludes the three printed intersections.

The following source remark instead says that these constructions differ
only in the intersection of the Suzuki subgroups. Deleting the last index
from the displayed intersection repairs that particular order defect,
but does not remove the independent rigidity obstruction to the intended
full lattices. We therefore record both conclusions rather than treating
an index correction as a proof of existence.

\paragraph{Attribution and scope.}
The source is D.~Leemans, \emph{The residually weakly primitive
pre-geometries of the Suzuki simple groups}, Note di Matematica 20(1)
(2000/2001), 1--20. The subgroup and normalizer premises are classical.
The rigidity proofs reproduce the project's earlier Suzuki realization
and higher-subfield notes; the contribution here is their explicit
application to the first four and sixth higher-rank constructions and
the distinction between literal and intended formulas.
This is not a classification of all higher-rank Suzuki geometries,
all relative positions, or the remaining constructions. It gives no
new principal-refutation, open-conjecture, Lean or compression credit.
\end{document}
